\chapter{Corporate Finance Core}\label{ch:corporate-finance-core}

% Scope: time value of money, NPV, the discount rate as hurdle, the growing perpetuity.
% Canonical reference for the book's house style: flowing prose under headings, with
% worked numbers set off in `example` and sharp cautions in `admonition`.
% The growing perpetuity's marketing twin lives in \cref{ch:customer-economics}.

\section{The Growing Perpetuity}\label{sec:growing-perpetuity}
% complexity: 3

What is a stream of cash worth today if it grows at a constant rate \(g\) forever? The
question is not academic: this single model underlies discounted-cash-flow terminal value
and customer lifetime value alike, and getting it wrong propagates through every valuation
that depends on them. The intuition is the foundation of all of finance. A dollar next year
is worth less than a dollar today, because today's dollar can be invested; growth only
partly offsets that discount. When the discount rate beats the growth rate, an infinite
stream of growing cash still sums to a finite number.

Summing the discounted cash flows and recognising a geometric series with ratio
\(\frac{1+g}{1+r}<1\) collapses the infinite sum to a single ratio:
\begin{equation}\label{eq:gordon}
  PV \;=\; \sum_{t=1}^{\infty}\frac{D\,(1+g)^{t-1}}{(1+r)^{t}}
      \;=\; \frac{D}{\,r-g\,}.
\end{equation}
The denominator \((r-g)\) is the net drag: discount pressure minus the growth lift. The
series converges \emph{only} when \(r>g\). If \(g\ge r\) the denominator turns zero or
negative and \cref{eq:gordon} returns nonsense --- the single most common valuation error,
and one a sceptical reader should check first in any model they are handed.

\begin{example}
A subscriber cohort throws off \money{1250000} next year at a \qty{42}{\percent} contribution
margin, growing \qty{3}{\percent} per year, and the firm discounts at \qty{11}{\percent}. The
margin is what actually accrues to the firm, so it is the \(D\) in \cref{eq:gordon}; the
\(\qty{11}{\percent}\) is the opportunity cost of the capital tied up in serving the cohort:
\[
  PV \;=\; \frac{\money{1250000}\times 0.42}{0.11-0.03}\;\approx\;\money{6562500}.
\]
The cohort is worth roughly \money{6562500} today. Note how sensitive that figure is to the
gap \(r-g\): widen it to \qty{10}{\percent} and the value falls by a fifth, which is why the
growth assumption deserves more scrutiny than any other input.
\end{example}

\begin{admonition}{Warning}
Reusing the same high \(g\) for ten explicit forecast years \emph{and} the terminal value
double-counts growth. The terminal \(g\) must reflect only the steady state --- usually at
or below long-run GDP growth. A terminal rate pulled from a company's recent hypergrowth is
the fastest way to manufacture a valuation that cannot survive contact with reality.
\end{admonition}

Relax constant growth and the model becomes a two-stage or H-model; relax constant churn in
its marketing twin and customer value needs a shifted-beta-geometric survival model. The
clean formula is the \emph{floor} of the analysis, not the ceiling. The identical algebra
reappears as customer lifetime value in \cref{ch:customer-economics}: rename \(D\) the annual
customer margin and \(g\) its growth rate, and \cref{eq:gordon} values a customer exactly as
it values a cohort. \textcite{brealey2020} develop the corporate-finance side in full.

\section{Time Value of Money}\label{sec:time-value-of-money}
% complexity: 2

How much is a future cash payment worth today? Every capital-allocation decision --- whether
to invest, acquire, or repurchase --- rests on answering this before anything else. A dollar
in hand today can be reinvested at rate \(r\), compounding to \((1+r)^n\) dollars in \(n\)
years; running that logic backward, a future payment is worth less than its face value by
exactly the factor compounding would have produced.

If a sum \(\mathrm{FV}\) arrives in \(n\) years and the opportunity cost of capital is \(r\)
per period, its present value follows by discounting:
\begin{equation}\label{eq:pv}
  PV \;=\; \frac{\mathrm{FV}}{(1+r)^{n}}.
\end{equation}
The denominator \((1+r)^{n}\) is the future-value factor; its reciprocal is the
\emph{discount factor}. \Cref{eq:gordon} is nothing but repeated application of \cref{eq:pv}
across an infinite, growing stream. The formula assumes \(r\) is constant and known, cash
arrives on the stated date, and the investor can actually reinvest at \(r\); for short
horizons or near-riskless flows the error from these assumptions is small, but for
long-dated, illiquid cash flows it can dominate the answer.

\begin{example}
What is a \money{252} payment arriving in three years worth today, at the firm's
\qty{11}{\percent} discount rate? That rate is the opportunity cost of equity capital derived
in \cref{ch:risk-and-capital-structure}; applying \cref{eq:pv},
\[
  PV \;=\; \frac{\money{252}}{(1.11)^{3}}
      \;=\; \frac{\money{252}}{1.3676}
      \;\approx\; \money{184}.
\]
The \money{68} gap between face value and present value is the cost of waiting --- not
inflation, but forgone compounding. The decision this informs: a promise to pay \money{252}
in three years is worth accepting only if the alternative use of \money{184} today earns less
than \qty{11}{\percent}.
\end{example}

\begin{admonition}{Warning}
Discounting nominal cash flows at a real rate, or vice versa, is the classic error.
Consistency between the flow and the rate is the only invariant of the present-value formula:
nominal flows need a nominal \(r\); inflation-adjusted flows need a real \(r\). Mixing them
misstates value by roughly the inflation rate per year, compounding quietly across a
multi-year project.
\end{admonition}

For risky cash flows the ``right'' \(r\) is itself a model output, not a given --- see
\cref{eq:capm} in \cref{ch:risk-and-capital-structure}. The present-value formula is the
atomic unit of all valuation work in this book; \textcite{brealey2020} derive it from first
principles and extend it to bonds, annuities, and the growing-perpetuity case of
\cref{eq:gordon}.

\section{Net Present Value}\label{sec:net-present-value}
% complexity: 3

Should the firm make an investment? Net present value converts an entire cash-flow schedule
--- outlay today, returns over time --- into a single number whose sign settles the question.
Positive means the investment earns more than the hurdle rate; negative means it destroys
value at that rate. The underlying question is plain: how much richer does this make the
investor, measured in today's dollars, after fully charging for the capital deployed?

Let \(\mathrm{FCF}_t\) be the free cash flow in period \(t\), \(I_0\) the initial outlay at
\(t=0\), and \(r\) the discount rate:
\begin{equation}\label{eq:npv}
  \mathrm{NPV} \;=\; \sum_{t=1}^{T}\frac{\mathrm{FCF}_t}{(1+r)^{t}} \;-\; I_0.
\end{equation}
Each term \(\mathrm{FCF}_t/(1+r)^t\) is \cref{eq:pv} applied to period \(t\). The discount
rate \(r\) is the firm's weighted average cost of capital, derived in
\cref{ch:risk-and-capital-structure} as \cref{eq:wacc}: it captures both the cost of equity
(\cref{eq:capm}) and the tax-shielded cost of debt, weighted by capital structure. The
formula assumes cash flows are known, \(r\) is constant, and the project carries no embedded
optionality --- no ability to expand, abandon, or defer --- so it can misprice investments
whose chief value is the flexibility they create.

\begin{example}
Does acquiring a customer at \money{600} create value at the firm's \qty{11}{\percent} hurdle
rate? Treat the \money{600} acquisition cost as the outlay \(I_0\). The customer's annual
contribution margin of \money{252} grows at \qty{3}{\percent}, so the inflow stream is a
growing perpetuity, valued by \cref{eq:gordon}:
\[
  \sum_{t=1}^{\infty}\frac{\mathrm{FCF}_t}{(1.11)^t}
  \;=\; \frac{\money{252}}{0.11-0.03}
  \;=\; \money{3150}.
\]
Applying \cref{eq:npv},
\[
  \mathrm{NPV} \;=\; \money{3150} \;-\; \money{600} \;=\; \money{2550}.
\]
Every dollar of acquisition cost returns \money{5.25} of present value --- a surplus far above
the hurdle. The implication is that the binding constraint here is not capital cost but
customer supply: the firm should raise its acquisition budget until the marginal cost of a
customer approaches the \money{3150} that customer is worth.
\end{example}

\begin{admonition}{Warning}
Using a too-low discount rate to manufacture a positive NPV is the cardinal sin of investment
appraisal. A project that looks attractive at \qty{6}{\percent} but is negative at the correct
\qty{11}{\percent} does not create value --- it destroys it slowly, until the true cost of
capital reasserts itself through refinancing pressure or equity dilution. Anchoring the hurdle
rate to the actual cost of capital is the discipline \cref{eq:npv} exists to enforce.
\end{admonition}

NPV is the corporate analogue of the customer-economics decision in
\cref{ch:customer-economics}: both ask whether a present outlay --- \(I_0\) or acquisition
cost --- is justified by the discounted stream it buys. The equivalence is not a coincidence;
it is the unifying logic of the value-creation spine. \textcite{brealey2020} prove NPV's
consistency with shareholder-value maximisation, and \textcite{ross2022} work through its
comparison with the internal rate of return and the conditions under which the two diverge.
